\documentclass[9pt,a4paper]{extarticle}

\usepackage[a4paper,top=0.62cm,bottom=0.62cm,left=0.95cm,right=0.95cm]{geometry}
\usepackage[T1]{fontenc}
\usepackage[utf8]{inputenc}
\usepackage{mathptmx}
\usepackage{amsmath}
\usepackage{array}
\usepackage{booktabs}
\usepackage{enumitem}
\usepackage{tabularx}
\usepackage{titlesec}
\usepackage{xcolor}
\usepackage{hyperref}
\usepackage{microtype}

\definecolor{ink}{HTML}{17324D}
\definecolor{accent}{HTML}{0E7490}
\definecolor{rulegray}{HTML}{B8C4CE}
\definecolor{sectionbg}{HTML}{EAF3F6}
\definecolor{muted}{HTML}{536575}

\hypersetup{
    colorlinks=true,
    linkcolor=accent,
    urlcolor=accent
}

\pagenumbering{gobble}
\setlength{\parindent}{0pt}
\setlength{\parskip}{0pt}
\setlength{\tabcolsep}{3.5pt}
\renewcommand{\arraystretch}{0.92}
\setlist[itemize]{leftmargin=1.2em,itemsep=1pt,topsep=1pt,parsep=0pt}
\setlist[enumerate]{leftmargin=1.35em,itemsep=1.5pt,topsep=1pt,parsep=0pt}
\setlength{\emergencystretch}{1.5em}
\raggedbottom

\titleformat{\section}
  {\bfseries\color{ink}\fontsize{9.2pt}{10pt}\selectfont}
  {}{0pt}
  {}
\titlespacing*{\section}{0pt}{2pt}{1pt}

\newcolumntype{L}{>{\raggedright\arraybackslash\bfseries\color{ink}\fontsize{7.8pt}{8.5pt}\selectfont}p{2.65cm}}
\newcolumntype{Y}{>{\raggedright\arraybackslash}X}
\newcommand{\proposalTable}[1]{%
    \renewcommand{\arraystretch}{0.94}%
    \begin{tabularx}{\textwidth}{@{}L Y@{}}
    \toprule[0.7pt]
    \textcolor{muted}{Field} & \textcolor{muted}{Content}\\
    \midrule
    #1
    \bottomrule[0.7pt]
    \end{tabularx}%
}

\begin{document}
\fontsize{8.4pt}{9.6pt}\selectfont

\begin{center}
    {\color{ink}\fontsize{12.5pt}{14pt}\selectfont\bfseries
    Deep Learning --- UAS Session 16: Group Project Proposal}\\[-1pt]
    {\color{accent}\fontsize{8.7pt}{10pt}\selectfont\bfseries
    Group Project Proposal Template --- Track A: Vision Classification with Transfer Learning}
\end{center}
\vspace{1pt}
\hrule height 0.8pt
\vspace{2pt}

\section*{1. Identification}
\proposalTable{%
    Group number & Kelompok 5\\
    Members & Tita Noviana, Titanio Yudista, Suci Fransisca Sisilia R., Fajar Dwiharjo\\
    Track & A (Vision classification with transfer learning)\\
    Proposed title & Evaluasi \textit{From-Scratch CNN} dan \textit{Fine-Tuned EfficientNet-B0} untuk \textit{Fine-Grained Classification} Motif Batik Semarang pada Kondisi Deformasi Kain\\
}

\section*{2. Problem Statement}
\begin{itemize}
    \item \textbf{Supported decision/prediction:} Mengklasifikasikan 10 motif batik Semarang untuk kurasi museum, pengarsipan warisan budaya, dan verifikasi katalog \textit{e-commerce}.
    \item \textbf{End users and failure cost:} Kurator, peneliti tekstil, dan pengelola \textit{platform} kriya. Prediksi keliru dapat mencemari arsip, salah memberi label produk, dan menurunkan kepercayaan pengguna.
    \item \textbf{Appropriateness of deep learning:} Citra menghadapi deformasi \textit{non-rigid}, lipatan, perubahan skala dan pencahayaan, serta kemiripan \textit{inter-class} pada \textit{isen-isen}; pola ini sulit ditangkap fitur manual seperti SIFT/HOG.
\end{itemize}

\section*{3. Data}
\proposalTable{%
    Dataset name & Semarang Batik Dataset\\
    Source URL & \url{https://www.kaggle.com/datasets/edywinarno/semarang-batik-dataset} (Winarno et al., ACM ICACS 2017)\\
    Licence or terms of use & Lisensi penggunaan publik dan akademik sesuai ketentuan Kaggle (\textit{Academic Research Use})\\
    Size and number of classes & $\sim 3.000$ citra RGB ($224\times224$) dalam 10 kelas motif Semarang\\
    Class balance & Relatif seimbang ($\sim 300$ citra/kelas); diverifikasi melalui \textit{EDA}\\
    Known missing or corrupted data & Skrip Python (PIL/OpenCV) memfilter \textit{corrupted headers} dan \textit{invalid color channels} sebelum \textit{splitting}\\
    Split procedure (train / validation / test) & \textit{Stratified split} 70\% / 15\% / 15\% dengan \textit{fixed random seed} (\texttt{seed=42})\\
    Leakage risks and how to check them & \textit{Exact/near-duplicates} diperiksa dengan \textit{pHash} sebelum \textit{splitting}; \textit{augmentation} hanya \textit{on-the-fly} pada \textit{training set}.\\
}

\section*{4. Baseline and Metric}
\proposalTable{%
    Baseline method & Arsitektur \textit{Custom 4-Block CNN} (\texttt{Conv2D + BatchNorm + ReLU + MaxPool + AdaptiveAvgPool + Linear head}) yang dilatih \textit{from-scratch}\\
    Why this baseline is a fair comparison & \textit{Data split}, resolusi, \textit{optimizer} (\texttt{AdamW}), dan \textit{epoch budget} identik; perbedaan utama hanya \textit{pretrained representations}.\\
    Primary metric & \textit{Macro F1-Score} untuk menilai 10 kelas secara seimbang\\
    Secondary metric & \textit{Top-1 Accuracy}, \textit{Top-3 Accuracy}, serta \textit{training time/cost per run}\\
    Success threshold & Model \textit{transfer learning} mencapai $\ge 85\%$ \textit{Macro F1-Score} dan mengungguli \textit{baseline CNN from-scratch} minimal $+10\%$ \textit{Macro F1} pada \textit{test set}\\
    Cost of false positive vs false negative & \textit{False Positive} mencemari arsip; \textit{False Negative} membuat motif asli terlewat dan memerlukan verifikasi manual.\\
}

\section*{5. Planned Comparison and Ablation}
\begin{enumerate}
    \item \textbf{Model comparison:} \textit{Custom 4-Layer CNN (From-Scratch)} vs \textit{Pretrained EfficientNet-B0 / ResNet-50}. \textit{Hierarchical features} ImageNet diharapkan lebih baik menangkap tekstur \textit{fine-grained}.
    \item \textbf{Ablation:} Membandingkan \textit{Geometric \& Photometric Data Augmentation (On vs Off)}---rotasi, \textit{affine transform}, dan \textit{color jitter}---dengan arsitektur dan \textit{hyperparameter} tetap.
    \item \textbf{Training implementation:} \textit{Custom PyTorch loop} untuk CNN \textit{from-scratch} dan \textit{two-stage fine-tuning} EfficientNet-B0: \textit{linear probing} lalu \textit{unfreezing top blocks} dengan \texttt{CosineAnnealingLR} dan \texttt{AdamW}.
\end{enumerate}

\section*{6. Compute Plan}
\noindent\textbf{Hardware:} Google Colab Free Tier GPU Tesla T4 (16GB VRAM). \textbf{Time:} $\sim 8$--$12$ menit per \textit{run} (30 \textit{epochs}, batch size 32, resolusi $224\times224$). \textbf{Runs:} 12 \textit{runs} (3 konfigurasi utama $\times$ 3 \textit{random seeds} + 3 \textit{sanity checks}). \textbf{Storage:} $\sim 600$ MB untuk \textit{checkpoints}, \textit{logs}, dan \textit{confusion matrix}. Rencana sesuai \textit{RPS compute policy}; jika terbatas, \textit{ablation} dipangkas menjadi 1 \textit{fixed seed} (\texttt{seed=42}).

\section*{7. AI Use Plan}
\noindent\textbf{Assistant tools:} Google Gemini dan GitHub Copilot. \textbf{Allowed:} \textit{debugging data loader}, kode \textit{boilerplate plotting}, dan perapian LaTeX. \textbf{Off-limits:} angka metrik fiktif atau analisis kausal \textit{error analysis} tanpa inspeksi visual citra kegagalan \textit{test set}. \textbf{Verification:} seluruh skrip dijalankan deterministik melalui satu \textit{single command} ber-\textit{fixed seed}; metrik divalidasi terhadap \textit{raw evaluation log}.

\end{document}
